\documentclass[10pt,a4paper]{amsart}
\usepackage[margin=19mm]{geometry}
\usepackage{amsmath,amssymb,mathtools}
\usepackage[scr=rsfs,cal=euler]{mathalfa}
\usepackage{xfrac}
\usepackage[unicode,hidelinks,bookmarks=false]{hyperref}
\newcommand{\ie}{i.e.\ }
\newcommand{\cf}{cf.\ }
\newcommand{\resp}{respectively\ }
\newcommand{\Subsection}{\subsection}
\providecommand{\nobreakdash}{\nobreak\hbox{-}}
\setlength{\emergencystretch}{2em}
\multlinegap0pt
\usepackage{mathtools}
\usepackage{enumerate}
\usepackage{tikz}
\usepackage[normalem]{ulem}
\usepackage{csquotes}
\newtheorem{lemma}{Lemma}
\newcommand\DD{{\mathbb D}}
\newcommand\TT{{\mathbb T}}
\title{Riesz $\alpha$-capacity of Cantor sets and cyclicity in Dirichlet-type spaces}
\author{Dimitrios Vavitsas, Jujie Wu, Konstantinos Zarvalis}
\date{Source: arXiv:2604.10324v1 (CC BY 4.0)}
\begin{document}
\maketitle

\noindent\emph{Source and licence.} Riesz $\alpha$-capacity of Cantor sets and cyclicity in Dirichlet-type spaces. Version arXiv:2604.10324v1 (CC BY 4.0), distributed by arXiv under the Creative Commons Attribution 4.0 International licence, http://creativecommons.org/licenses/by/4.0/, as recorded in the arXiv licence metadata for this exact version and shown on the abstract page https://arxiv.org/abs/2604.10324v1. Full licence notice: \texttt{licenses/arxiv-2604.10324-CC-BY-4.0.txt}.\par\smallskip

\noindent\textit{Editorial scope.} Counted unit: the article's lemma estimation (source0.tex lines 415--421). The article's complete original proof of it is delivered with it (source0.tex lines 422--464, one \texttt{proof} environment of 43 lines). No mathematics is written, repaired or re-derived: the statement and the proof are the article's own lines, in the article's own notation, and no retained line is substituted or re-flowed.  No further source-local context is needed: every cross-reference of every retained line resolves inside the two delivered spans.  Nothing was omitted from any retained span, so the package carries no omission record.  No retained line contains a citation, so the wrapper carries no bibliography and no bibliography entry.  No retained line contains a figure environment, an image-inclusion command or drawing code, so no image is delivered and the package declares no asset dependency.  Editorial formatting only, in addition: an AMS \texttt{amsart} preamble that restates the article's own declarations and macros used by the retained lines, namely the environments \texttt{lemma} on separate counters, together with the standard packages those lines need.  The retained environments print in this document as Lemma 1 in order of appearance, whereas the article prints the same items with the numbering of its own full document.  The complete native source of the article as distributed in the arXiv e-print for this exact version is bundled unchanged as \texttt{sources/198264/source0.tex}.
\begin{lemma}\label{Real estimation}
Fix $\alpha_1\in (0,1)$. There exists a constant $C\ge1$ depending solely on the choice of $\alpha_1$ such that
\begin{equation*}
\frac{1}{|1-z\overline{\eta}|^{1-\alpha}}\leq C \text{ }\mathrm{Re}\Big(\frac{1}{(1-z\overline{\eta})^{1-\alpha}}\Big),  
\end{equation*} 
for all $\alpha\in [\alpha_1,1]$, $z\in\DD$ and $\eta\in\TT$.
\end{lemma}

\begin{proof}
Set $z=re^{i\varphi}\in\DD$ and $\eta=e^{i\theta}\in \TT$, where $\varphi,\theta\in [0,2\pi)$ and $r\in[0,1)$. Note that
\begin{equation*}
1-z\overline{\eta}=1-re^{i(\varphi-\theta)}=1-r\cos(\varphi-\theta)-ir\sin(\varphi-\theta),
\end{equation*}
and thus
\begin{equation*}
  \mathrm{Re}(1-z\overline{\eta})=1-r\cos(\varphi-\theta)>0.  
\end{equation*}
Thus, we may write
\begin{equation}\label{arg}
\arg{(1-z\overline{\eta}})=\arctan\frac{\mathrm{Im}{(1-z\overline{\eta}})}{\mathrm{Re}{(1-z\overline{\eta}})}=\arctan\frac{-r\sin(\varphi-\theta)}{1-r\cos(\varphi-\theta)}\in\left(-\frac{\pi}{2},\frac{\pi}{2}\right).
\end{equation}
In case $\alpha=1$, the statement of the lemma is trivially true for any $C\ge1$. Next, pick $\alpha\in [\alpha_1,1)$. One may see that
\begin{equation}\label{Re}
\mathrm{Re}\Big(\frac{1}{(1-z\overline{\eta})^{1-\alpha}}\Big)=\frac{\mathrm{Re}(e^{-i(1-\alpha)\arg(1-z\overline{\eta})})}{|1-z\overline{\eta}|^{1-\alpha}}=\frac{\cos((1-\alpha)\arg(1-z\overline{\eta}))}{|1-z\overline{\eta}|^{1-\alpha}}.
\end{equation}
We wish to prove the existence of an absolute constant $C\ge1$ such that 
\begin{equation*}
   1\leq C \text{ }|1-z\overline{\eta}|^{1-\alpha}\mathrm{Re}\Big(\frac{1}{(1-z\overline{\eta})^{1-\alpha}}\Big), 
\end{equation*}
which, via \eqref{arg} and \eqref{Re} is equivalent to
\begin{equation*}
1\leq C\cos{\Big((1-\alpha) \arctan\frac{-r\sin(\varphi-\theta)}{1-r\cos(\varphi-\theta)}\Big)},
\end{equation*}
and $C\ge1$ does not depend on $r\in [0,1)$, $\varphi,\theta\in [0,2\pi)$, and $\alpha\in [\alpha_1,1)$. In particular, such a uniform constant exists if
\begin{equation}
\inf_{\alpha,r, \varphi,\theta}{\cos{\Big((1-\alpha) \arctan\frac{-r\sin(\varphi-\theta)}{1-r\cos(\varphi-\theta)}\Big)}}>0,
\end{equation}
where $ \alpha,r, \varphi,\theta$ lie in the sets mentioned above. Indeed, since the cosine and arctangent functions are even and odd, respectively, we have
\begin{equation*}
\cos{\Big((1-\alpha) \arctan\frac{-r\sin(\varphi-\theta)}{1-r\cos(\varphi-\theta)}\Big)}=\cos{\Big((1-\alpha) \arctan\frac{r|\sin(\varphi-\theta)|}{1-r\cos(\varphi-\theta)}\Big)}.
\end{equation*}
Furthermore,
\begin{equation*}
    (1-\alpha)\arctan\frac{r|\sin(\varphi-\theta)|}{1-r\cos(\varphi-\theta)}\in [0,(1-\alpha)\pi/2),
\end{equation*}
and hence,  by the monotonicity of $\cos x$ and $\arctan x$ on $[0,\pi/2)$, we obtain
\begin{equation*}
\cos{\Big((1-\alpha) \arctan\frac{r|\sin(\varphi-\theta)|}{1-r\cos(\varphi-\theta)}\Big)}\geq \cos\left((1-\alpha)\frac{\pi}{2}\right)\geq \cos\left((1-\alpha_1)\frac{\pi}{2}\right)>0,
\end{equation*}
for all $r\in [0,1)$, $\varphi,\theta\in [0,2\pi)$, and $\alpha\in [\alpha_1,1)$. This proves the assertion with $C=(\cos((1-\alpha_1)\pi/2))^{-1}>1$.
\end{proof}

\end{document}
